\documentclass[11pt]{article}
\usepackage{amsmath,amssymb,amsthm,mathtools}
\usepackage{enumitem}
\usepackage[margin=1in]{geometry}

\newtheorem{definition}{Definition}
\newtheorem{theorem}{Theorem}
\newtheorem{nogo}{No-Go Theorem}

\title{Step 162: Calkin Algebra Declaration and Source-Theorem Extraction Audit}
\author{Six Birds / RH Membrane Working Note}
\date{}

\begin{document}
\maketitle

\section{Inherited state}

The track remains \emph{adequacy-oriented}.  The active residual is
\[
  \Xi^{\rm BC}_{\ell,a}
  =
  \mathfrak B_{\ell,a}^{*}\Pi_{Y_a}\mathfrak B_{\ell,a},
\]
and the active operator is
\[
  C_\ell P_\eta,
  \qquad
  C_\ell=(I-P_\infty)M_{m_\ell}P_\infty,
  \qquad
  m_\ell(s)=e^{-\ell(1/2-s)}.
\]
The Step 161 verdict was that no local Calkin-faithful boundary-symbol theorem
had been extracted.  Step 162 refines that verdict by declaring the exact
Calkin algebra and splitting the missing bridge theorem into source-testable
sub-obligations.

\section{Typed Calkin algebra}

\begin{definition}[Pulled-evaluator Calkin algebra]
Let \(H\) be the Burnol/Sonine pulled-evaluator Hilbert context containing
\[
  H_\eta=
  \overline{\operatorname{span}\{J_a^*Y^a_{\rho,k}\}}.
\]
Let
\[
  \mathcal A_\eta
  =
  C^*(P_\infty,M_{m_\ell},P_\eta,I)
  \subseteq \mathcal B(H),
\]
where \(P_\eta\) is the projection onto \(H_\eta\).  The compact ideal in this
context is
\[
  \mathcal K_\eta=\mathcal A_\eta\cap\mathcal K(H).
\]
The quotient map is
\[
  q_\eta:\mathcal A_\eta\to\mathcal A_\eta/\mathcal K_\eta.
\]
\end{definition}

The active quotient object is
\[
  q_\eta(C_\ell P_\eta)
  \in
  \mathcal A_\eta/\mathcal K_\eta.
\]
This is the object that a boundary-symbol theorem must classify.  A theorem on
an ambient Hardy algebra, a hard-support algebra, or a finite section is not a
theorem about \(q_\eta(C_\ell P_\eta)\) unless it is connected to this quotient
by an accepted bridge.

\section{Split bridge theorem}

Step 162 rewrites the missing Step 161 bridge as five sub-obligations:
\begin{enumerate}[label=(G\arabic*)]
  \item \textbf{Algebra gate.} The source theorem is stated on
  \(\mathcal A_\eta/\mathcal K_\eta\), or supplies an accepted bridge into it.
  \item \textbf{Symbol gate.} There is a boundary symbol
  \[
    \sigma_B:\mathcal A_\eta/\mathcal K_\eta\to\mathcal S_B
  \]
  faithful on \(q_\eta(C_\ell P_\eta)\).
  \item \textbf{Normal-form gate.} The active operator admits
  \[
    C_\ell=U_\ell W_\ell+K_\ell,
    \qquad
    K_\ell\in\mathcal K(H).
  \]
  \item \textbf{Lower-faithfulness gate.} \(U_\ell\) is bounded below on the
  boundary range relevant to \(H_\eta\).
  \item \textbf{Compact-ideal gate.} Compactness of \(W_\ell P_\eta\) is
  equivalent to compactness of \(C_\ell P_\eta\) in the declared quotient.
\end{enumerate}

These gates are independent.  Supplying only one of them does not accept the
bridge.

\section{Source extraction audit}

The audited sources split as follows.

\begin{center}
\begin{tabular}{llll}
\hline
Source & Provides & Missing & Verdict\\
\hline
Burnol & carrier and zero-evaluator criterion & Calkin symbol and normal form & support only\\
Steps 104--105 & raw Sonin Calkin obstruction & bridge to \(H_\eta\) and repair symbol & structural input\\
Steps 114--115 & co-Poisson exactness criterion & shifted zeta factorization & open route\\
Step 154 & commutator formula \(C_\ell\eta\) & annihilation or compactness theorem & open route\\
Steps 160--161 & conditional bridge theorem and source audit & bridge existence & conditional\\
Hardy/hard-support shadows & candidate symbol intuition & typed carrier bridge & public shadow\\
\hline
\end{tabular}
\end{center}

No source passes all five gates.  The status is therefore
\[
  \boxed{\texttt{split\_external\_theorem}.}
\]

\section{Retained framework results}

\begin{nogo}[Public-shadow discipline]
A Hardy/Hankel or hard-support symbol cannot decide \(q_\eta(C_\ell P_\eta)\)
without an accepted bridge into \(\mathcal A_\eta/\mathcal K_\eta\).
\end{nogo}

\begin{nogo}[Finite-window Calkin blindness]
Finite singular-value data can diagnose likely structure but cannot certify the
completed Calkin class of \(C_\ell P_\eta\).
\end{nogo}

\begin{theorem}[Conditional bridge normal form]
If a source supplies gates G1--G5 above, then the Step 160 equivalence applies:
\[
  C_\ell P_\eta\in\mathcal K(H)
  \quad\Longleftrightarrow\quad
  W_\ell P_\eta\in\mathcal K(H).
\]
\end{theorem}

\section{Verdict and next leaf}

Step 162 makes the external theorem proof-ready but does not prove it.  The
missing theorem is now split into:
\[
  \boxed{
  \mathcal A_\eta/\mathcal K_\eta
  \;+\;
  \sigma_B
  \;+\;
  C_\ell=U_\ell W_\ell+K_\ell
  \;+\;
  U_\ell\text{ lower-faithful}
  \;+\;
  K_\ell\in\mathcal K(H).
  }
\]
The next step should choose one branch:
\begin{enumerate}[label=(\alph*)]
  \item prove the split theorem from CCM/Sonin operator theory;
  \item build a pulled-evaluator Weyl sequence in \(H_\eta\);
  \item prove shifted co-Poisson factorization; or
  \item record a scoped residual theorem carrying \(\Xi^{\rm BC}\).
\end{enumerate}

\end{document}
